\chapter{Trabajos relacionados}
\label{cap:trabajos-relacionados}

Este cap\'itulo integra la revisi\'on de trabajos relacionados realizada por
Zaida (bloque requisitos/usabilidad/dominio veterinario), Jaime (bloque
arquitectura/seguridad/rendimiento/calidad) y Fred (bloque arquitectura de
datos/rendimiento), consolidada y cerrada con un protocolo inspirado en
PRISMA~2020~\cite{page2021prisma} en \texttt{docs/checklists/prisma2020.md}.
PRISMA fue dise\~nado para revisiones sistem\'aticas de intervenciones
sanitarias; aqu\'i se adapta expl\'icitamente como marco de \textbf{rigor y
transparencia} para una revisi\'on narrativa de trabajos relacionados de
ingenier\'ia de software, no como un metaan\'alisis cl\'inico --- varios
\'items orientados a estad\'istica de efectos cl\'inicos se marcan
\emph{No aplica} en el checklist original en vez de forzarlos.

\section{Criterios de elegibilidad y fuentes}

\textbf{Inclusi\'on:} estudio con DOI o enlace verificable; publicado en
fuente acad\'emica o en una revista/\emph{proceedings} con revisi\'on por
pares (o norma/documento oficial de industria sin DOI, como ISO u OWASP);
relevante a al menos uno de los tres bloques tem\'aticos (requisitos,
usabilidad o dominio veterinario --- Zaida; arquitectura, seguridad,
rendimiento o calidad --- Jaime; arquitectura de datos/rendimiento ---
Fred). \textbf{Exclusi\'on:} blogs comerciales sin revisi\'on por pares,
p\'aginas de venta de software, trabajos sin autor o fecha verificable.
\textbf{Fuentes consultadas:} IEEE Xplore, ACM Digital Library, Scopus y
Google Scholar (bloque Zaida, b\'usqueda 2026-08-17); auditor\'ia de la
bibliograf\'ia ya existente m\'as verificaci\'on dirigida contra Crossref y
p\'aginas oficiales de editorial/conferencia/norma (bloque Jaime, misma
fecha); verificaci\'on dirigida de DOI reales, uno a uno, con un filtro
temporal expl\'icito de 2020--2026 (bloque Fred, 2026-08-17). El proceso de
selecci\'on fue independiente por bloque (cada autor filtr\'o su propio
tema), sin doble revisor cruzado por estudio --- limitaci\'on declarada
expl\'icitamente, no oculta.

De un total de 15 candidatos revisados por Zaida (12 excluidos por no
cumplir el criterio de elegibilidad), una verificaci\'on dirigida de 4
candidatos por Jaime, y 13 candidatos por DOI verificados uno a uno por
Fred, el conjunto final incluido en la \textbf{s\'intesis de trabajos
relacionados de este cap\'itulo} suma \textbf{10 estudios} (Zaida~3 +
Jaime~4 + Fred~3), cumpliendo la meta original de $\geq$8--9. Esos
\textbf{10} son espec\'ificamente los estudios incluidos en esta
s\'intesis narrativa --- no el total de referencias retenidas por los
tres bloques en todo el informe.

\textbf{Trazabilidad completa del bloque Fred (corregida 2026-09-04, tras
auditor\'ia cruzada de Jaime y verificaci\'on de Fred).} La etapa de
\emph{identificaci\'on} parti\'o de 13 candidatos por DOI (rango
2020--2026). La etapa de \emph{verificaci\'on} (17-08-2026) descart\'o 2
(Yang et al., 2018, ICSE --- fuera del rango temporal; un art\'iculo de
TURCOMAT --- inaccesible, HTTP~403), dejando \textbf{11 referencias
elegibles/retenidas}, las 11 con DOI verificado (ver
\texttt{docs/investigacion/handoff-fred-referencias.bib}). Sobre esas 11,
se aplic\'o una \textbf{clasificaci\'on editorial para este informe ---
no un paso de exclusi\'on PRISMA}: \textbf{3} son trabajos relacionados
de dominio/arquitectura y se incorporan a la s\'intesis de este cap\'itulo
(Bucao et al., Llaneta et al., Zmaranda et al. ---
secci\'on~\ref{sec:sintesis-trabajos-relacionados}); las otras
\textbf{8 no fueron excluidas en ning\'un momento}: son referencias de
apoyo metodol\'ogico/rendimiento citadas en otros cap\'itulos del informe
(arquitectura de acceso a datos, cach\'e, despliegue contenerizado ---
\texttt{nowicki2022comparative}, \texttt{zuchnik2021comparative},
\texttt{siebyla2020impact}, \texttt{dhalla2021performance},
\texttt{godinho2024performance}, \texttt{zulfa2020caching},
\texttt{yang2021large}, \texttt{ghatrehsamani2020art}), fuera del alcance
de la s\'intesis narrativa de trabajos relacionados pero citadas con
funci\'on propia en el informe. Los tres bloques est\'an integrados en el
checklist PRISMA (\texttt{docs/checklists/prisma2020.md}) y en el
diagrama de flujo que se muestra en la figura~\ref{fig:prisma-flow}.

% EVIDENCIA-COMPARTIDA-PRISMA
% Ruta esperada: figuras/zaida/06-prisma-flow-diagram.png
\evidencia{figuras/zaida/06-prisma-flow-diagram.png}%
{Diagrama de flujo PRISMA de identificaci\'on y selecci\'on de estudios (regenerado 2026-09-04, trazabilidad del bloque Fred corregida y verificada por Fred). Bloque Zaida completamente cuantificado (15 identificados $\to$ 12 excluidos $\to$ 3 incluidos); bloque Fred completamente cuantificado (identificaci\'on: 13 candidatos $\to$ verificaci\'on, 17-08-2026: 2 descartados $\to$ 11 elegibles/retenidos $\to$ clasificaci\'on editorial para este informe, no un paso de exclusi\'on PRISMA: 3 trabajos relacionados + 8 de apoyo metodol\'ogico/rendimiento, estos \'ultimos citados en otros cap\'itulos, no en la s\'intesis de n=10); bloque Jaime con el conteo final de incluidos (4), sin registro de candidatos excluidos (\'unico gap que persiste, se\~nalado con l\'inea punteada).}%
{fig:prisma-flow}

\section{Cadena de búsqueda, fuentes y ventana temporal por bloque}
\label{sec:cadena-busqueda-trabajos-relacionados}

Los tres bloques usaron estrategias de búsqueda distintas entre sí ---
diferencia declarada explícitamente aquí, no oculta detrás del checklist.
A continuación, la cadena exacta (cuando existió), las bases/fuentes y la
ventana temporal aplicada por cada uno, tal como quedaron registradas el
2026-08-17 en \texttt{docs/informe/borradores/zaida/estudios-primarios-zaida.md},
\texttt{docs/informe/borradores/jaime/trabajos-relacionados.md} y
\texttt{docs/investigacion/handoff-fred-trabajos-relacionados.md}
respectivamente (fuente primaria de cada bloque; el detalle completo,
incluida la tabla de exclusiones, sigue disponible en
\texttt{docs/checklists/prisma2020.md}, sección 4, ítems 6--7).

\textbf{Bloque Zaida (requisitos, usabilidad, dominio veterinario).}
Bases consultadas: IEEE Xplore, ACM Digital Library, Scopus, Google
Scholar. Fecha de búsqueda: 2026-08-17. \textbf{Sin ventana temporal
declarada} (no se aplicó filtro de año como criterio formal). Tres
cadenas de búsqueda booleanas exactas, una por sub-tema:

\begin{quote}
\texttt{"requirements elicitation" AND "systematic literature review"} \\
\texttt{"system usability scale" AND web application} \\
\texttt{veterinary management system web application case study}
\end{quote}

Resultado: 15 candidatos revisados, 12 excluidos por no cumplir el
criterio de elegibilidad de la sección anterior, 3 incluidos.

\textbf{Bloque Jaime (arquitectura, seguridad, rendimiento, calidad).}
\textbf{No usó una cadena booleana única.} Se auditó primero la
bibliografía ya existente en \texttt{docs/informe/referencias.bib} (15
entradas, todas normas/documentación oficial, ninguna un trabajo
académico de \emph{related work}) y, al no encontrar ahí trabajos
académicos pertinentes, se hizo verificación dirigida por tema
(metodología ágil de gestión veterinaria, pruebas de rendimiento web,
vulnerabilidades JWT, aplicación práctica de ISO/IEC~25010) contra
Crossref (\texttt{api.crossref.org}) y páginas oficiales de
editorial/conferencia. Fecha: 2026-08-17. \textbf{Sin ventana temporal
declarada.} Esta es una diferencia metodológica real frente al bloque de
Zaida, no un vacío de documentación: se trató de verificación de
referencias ya identificadas por relevancia temática directa, no de una
búsqueda exploratoria amplia que requiriera una cadena. Resultado: 4
estudios incluidos; el conteo de candidatos evaluados y excluidos antes
de esos 4 no fue registrado por Jaime (limitación declarada en la
sección~\ref{sec:limitaciones-trabajos-relacionados} y en el diagrama de
flujo, figura~\ref{fig:prisma-flow}).

\textbf{Bloque Fred (arquitectura de datos, rendimiento).} Tampoco usó
una cadena booleana: verificación dirigida de DOI reales, uno a uno,
abriendo cada \texttt{https://doi.org/<DOI>} y confirmando título,
autores y año contra Crossref/IEEE/ACM/IOP. Fecha: 2026-08-17.
\textbf{Única ventana temporal explícita de los tres bloques: 2020--2026}
(criterio propio de Fred, ``máx. 5--6 años''). Resultado: 13 candidatos
por DOI, 2 descartados (Yang et al., ICSE 2018, DOI
\texttt{10.1145/3180155.3180194} --- fuera del rango 2020--2026; un
artículo de TURCOMAT --- inaccesible, HTTP~403), 3 incluidos.

\section{Estudios incluidos}

\subsection{Pacheco, Garc\'ia \& Reyes (2018) --- T\'ecnicas de elicitaci\'on de requisitos}
Revisi\'on sistem\'atica sobre 140 estudios (1993--2015) que identifica las
entrevistas estructuradas como la t\'ecnica de elicitaci\'on de requisitos
con mayor evidencia emp\'irica de madurez y efectividad~\cite{pacheco2018requirements}.
El SRS de BIOPET (\texttt{docs/requisitos/SRS.md}) documenta requisitos ya
elicitados (historias de usuario, casos de uso, patr\'on ISO/IEC/IEEE~29148)
pero no declara expl\'icitamente qu\'e t\'ecnica se us\'o para llegar a
ellos; a diferencia de Pacheco et al., que comparan t\'ecnicas entre s\'i,
BIOPET aplica una t\'ecnica ya elegida y la documenta con trazabilidad
formal completa (SRS~$\to$~HU~$\to$~CU~$\to$~matriz), sin reportar esa
comparaci\'on.

\subsection{Bangor, Kortum \& Miller (2008) --- Benchmark emp\'irico del SUS}
Analiza aproximadamente diez a\~nos de datos SUS sobre cientos de
productos, estableciendo un puntaje medio de referencia de la industria de
68/100 (``sobre el promedio'') y 80{,}3 como umbral de
``excelente''~\cite{bangor2008sus}. Este estudio es la fuente del
benchmark contra el que se interpreta el puntaje SUS propio de BIOPET
(secci\'on~\ref{sec:sus-final}): BIOPET no construye un benchmark nuevo,
usa el ya existente para interpretar una medici\'on real sobre un sistema
concreto.

\subsection{Cede\~no Ochoa, Catuto Murillo \& Rodas-Silva (2021) --- Aplicaciones web para cl\'inicas veterinarias en Ecuador}
Caso de estudio sobre el desarrollo de una aplicaci\'on web para gesti\'on
de una cl\'inica veterinaria ecuatoriana, documentando la mejora
cualitativa de procesos administrativos al migrar de manejo manual a una
plataforma centralizada~\cite{cedeno2021veterinaria}. Es el trabajo con
contexto pa\'is m\'as cercano encontrado y verificado. A diferencia de
BIOPET, no reporta una evaluaci\'on emp\'irica multim\'etrica del sistema
resultante (sin cobertura de pruebas, rendimiento bajo carga, seguridad
auditada ni usabilidad medida con instrumento estandarizado).

\subsection{Rodr\'iguez, Llerena, Guevara, Baren \& Castro (2024) --- OSCRUM}
Estudio de caso sobre el desarrollo de SysVet, un sistema web de c\'odigo
abierto para gesti\'on veterinaria, usando la metodolog\'ia \'agil
OSCRUM~\cite{rodriguez2024oscrum}. Comparte con BIOPET el inter\'es en la
reproducibilidad del artefacto (c\'odigo abierto, proceso documentado),
pero documenta el \emph{proceso} de desarrollo \'agil, no una evaluaci\'on
emp\'irica multim\'etrica del \emph{producto} resultante.

\subsection{Putri, Hadi \& Ramdani (2017) --- Pruebas de rendimiento de un sistema web acad\'emico}
Caracteriza cuantitativamente el rendimiento de un sistema de admisi\'on
de estudiantes bajo carga, mediante \emph{benchmarking} t\'ecnico contra el
sistema real~\cite{putri2017perftesting}. Es metodol\'ogicamente an\'alogo a
la evaluaci\'on de rendimiento de BIOPET con k6
(secci\'on~\ref{sec:rendimiento-final}): ambos son estudios de
\emph{benchmarking} de un artefacto web real en un entorno acad\'emico, sin
dise\~no experimental controlado. A diferencia de BIOPET, Putri et al.
miden \'unicamente rendimiento, sin seguridad, usabilidad ni cobertura de
pruebas.

\subsection{Yang et al. (2026) --- Vulnerabilidades en implementaciones de JWT (NDSS)}
Auditor\'ia sistem\'atica de 43 implementaciones reales de JSON Web Token
en diez lenguajes, con una herramienta de \emph{fuzzing} dirigido
(JWTeemo), que descubre 31 vulnerabilidades previamente desconocidas (20
con CVE asignado)~\cite{yang2026tokentimebomb}. BIOPET usa JWT como
mecanismo central de autenticaci\'on (cap\'itulo~\ref{cap:seguridad}); este
trabajo es la referencia acad\'emica m\'as actual disponible sobre los
riesgos reales de esa misma tecnolog\'ia. A diferencia de Yang et al., que
auditan \emph{implementaciones/librer\'ias} gen\'ericas, BIOPET no reclama
haber descubierto vulnerabilidades JWT como contribuci\'on cient\'ifica:
aporta evidencia de que una aplicaci\'on concreta que usa JWT fue auditada
con m\'ultiples t\'ecnicas y no present\'o los patrones de mayor severidad
detectables por ellas, sin afirmar estar ``libre'' de la categor\'ia de
vulnerabilidades de \emph{fuzzing} dirigido a la librer\'ia subyacente
(\texttt{jjwt}), que queda expl\'icitamente fuera del alcance de la
auditor\'ia de BIOPET.

\subsection{Estdale \& Georgiadou (2018) --- Aplicaci\'on pr\'actica de ISO/IEC~25010}
Discute la dificultad pr\'actica de aplicar el modelo de calidad
ISO/IEC~25010 a un producto de software real, m\'as all\'a de la cita
te\'orica de la norma~\cite{estdale2018applying25010}. Es paralelo
metodol\'ogico directo al mapeo realizado en
\texttt{docs/arquitectura/ISO-25010.md} (secci\'on~\ref{sec:iso25010}),
donde se mapean las ocho caracter\'isticas de ISO/IEC~25010:2011 contra
evidencia real de BIOPET, declarando honestamente ``no evaluada
directamente'' cuando no existe evidencia. A diferencia del trabajo de
Estdale \& Georgiadou, que discute el modelo en general, BIOPET aporta un
mapeo concreto con rutas de evidencia verificables por cada
subcaracter\'istica.

\subsection{Bucao, Quinones, Abong, Penuela \& Sun (2023) --- VETelgeuse}
Sistema de gesti\'on web y m\'ovil para cl\'inicas veterinarias peque\~nas y
medianas, evaluado con el cuestionario de usabilidad USE sobre 30
usuarios reales~\cite{bucao2023vetelgeuse}. Confirma, junto con Cede\~no
Ochoa et al. y Llaneta et al., que el dominio de BIOPET (gesti\'on
veterinaria web) tiene evidencia acad\'emica real m\'as all\'a de un
\'unico estudio. A diferencia de BIOPET, su validaci\'on emp\'irica se
limita a usabilidad, sin cobertura de pruebas, rendimiento bajo carga ni
auditor\'ia de seguridad.

\subsection{Llaneta, Guelas, Labanan, Mercado \& Sasis (2022) --- TerraVet}
Propone un \emph{framework} web y m\'ovil para conectar due\~nos de
mascotas con cl\'inicas veterinarias, publicado en ICIST~2022 (ACM), con
una evaluaci\'on preliminar de la propuesta~\cite{llaneta2022terravet}. Es
paralelo de dominio y de arquitectura web+m\'ovil a BIOPET, pero no eval\'ua
rendimiento bajo carga ni despliegue en la nube, ambos reportados en este
informe (cap\'itulos~\ref{cap:pruebas-calidad} y \ref{cap:despliegue}).

\subsection{Zmaranda, Pop-Fele, Gyor\"odi, Gyor\"odi \& Pecherle (2020) --- CRUD con ORM vs. SQL directo}
Compara el tiempo de ejecuci\'on de operaciones CRUD usando mapeadores
objeto-relacionales (ORM) de .NET frente a acceso directo a
datos~\cite{zmaranda2020crud}. Respalda directamente la decisi\'on de
arquitectura de BIOPET de usar un \textbf{acceso h\'ibrido}: Spring Data
JPA para operaciones CRUD simples y SQL nativo en procedimientos
almacenados para consultas complejas y validaciones cruzadas
(secci\'on~\ref{sec:acceso-jpa-formal}). A diferencia de BIOPET, el estudio
de Zmaranda et al. se limita a una sola pila (.NET) y no eval\'ua
procedimientos almacenados ni agregaciones equivalentes a
\texttt{fn\_reporte\_dashboard} o \texttt{fn\_historial\_clinico\_mascota}.

\section{S\'intesis comparativa}
\label{sec:sintesis-trabajos-relacionados}

\begin{table}[H]
  \centering
  \tiny
  \caption{S\'intesis comparativa de los diez estudios incluidos frente a BIOPET.}
  \label{tab:sintesis-relacionados}
  \begin{tabular}{@{}p{2.4cm}p{2.2cm}p{3.0cm}p{4.8cm}@{}}
    \toprule
    Estudio & Dominio & Evaluaci\'on reportada & Diferencia principal frente a BIOPET \\
    \midrule
    Pacheco et al. (2018) & Elicitaci\'on de requisitos & Madurez/efectividad comparada de t\'ecnicas & BIOPET no compara t\'ecnicas; aplica una y la traza formalmente \\
    Bangor et al. (2008) & Usabilidad (SUS) & Benchmark de interpretaci\'on (68/80{,}3) & BIOPET usa el benchmark, no construye uno nuevo \\
    Cede\~no Ochoa et al. (2021) & Gesti\'on veterinaria (Ecuador) & Mejora cualitativa de procesos & Sin evaluaci\'on multim\'etrica; BIOPET s\'i \\
    Rodr\'iguez et al. (2024) --- OSCRUM & Gesti\'on veterinaria, c\'odigo abierto & Proceso \'agil documentado & Documenta proceso, no eval\'ua emp\'iricamente el producto \\
    Putri et al. (2017) & Sistema web acad\'emico & Rendimiento bajo carga & Solo rendimiento; sin seguridad, usabilidad ni cobertura \\
    Yang et al. (2026) --- NDSS & Seguridad JWT (multi-lenguaje) & Vulnerabilidades descubiertas & Audita librer\'ias gen\'ericas, no una aplicaci\'on completa \\
    Estdale \& Georgiadou (2018) & Calidad de software (ISO/IEC~25010) & Discusi\'on de evidenciabilidad & Gu\'ia general, no un mapeo concreto enlazado a evidencia \\
    Bucao et al. (2023) --- VETelgeuse & Gesti\'on veterinaria (Filipinas) & Usabilidad (USE, n=30) & Solo usabilidad; sin rendimiento, seguridad ni cobertura \\
    Llaneta et al. (2022) --- TerraVet & Gesti\'on veterinaria (marco web+m\'ovil) & Propuesta de \emph{framework}, evaluaci\'on preliminar & No eval\'ua rendimiento ni despliegue en la nube \\
    Zmaranda et al. (2020) & Acceso a datos (ORM vs. SQL directo) & Tiempos de operaciones CRUD & Limitado a .NET; no cubre procedimientos almacenados ni agregaciones \\
    \bottomrule
  \end{tabular}
\end{table}

Como resume la tabla~\ref{tab:sintesis-relacionados}, tomados
individualmente los diez estudios eval\'uan siempre \textbf{una
sola dimensi\'on a la vez}: elicitaci\'on de requisitos sin evaluar el
producto resultante, un benchmark de usabilidad sin aplicarlo a un sistema
concreto, el dominio veterinario sin evaluaci\'on multim\'etrica (tres
veces: Cede\~no Ochoa, Bucao, Llaneta), un proceso \'agil sin evaluaci\'on
emp\'irica del producto, rendimiento sin seguridad ni usabilidad, seguridad
JWT gen\'erica sin evaluar una aplicaci\'on completa, el modelo de calidad
discutido sin aplicarlo con evidencia enlazada a un caso real, o
rendimiento de acceso a datos limitado a una sola pila tecnol\'ogica.
BIOPET se posiciona frente a estos diez trabajos como una evaluaci\'on
\textbf{multim\'etrica} de un \'unico artefacto: cobertura de pruebas,
rendimiento, usabilidad, calidad web, y seguridad por tres \'angulos
distintos (revisi\'on manual OWASP, an\'alisis din\'amico ZAP, an\'alisis
est\'atico SpotBugs/Find Security Bugs), todas enmarcadas bajo
ISO/IEC~25010 como vocabulario com\'un. Con la incorporaci\'on del bloque
de Fred, el dominio veterinario exacto de BIOPET pasa de tener un solo
estudio de contraste (Cede\~no Ochoa et al.) a tres (sumando Bucao et al. y
Llaneta et al.), aunque ninguno de los tres combina el contexto
ecuatoriano con una evaluaci\'on multim\'etrica --- esa combinaci\'on
espec\'ifica sigue sin evidencia externa directa. No se afirma que BIOPET
sea metodol\'ogicamente superior a ninguno de estos trabajos: no existe una
comparaci\'on emp\'irica directa que lo demuestre; la diferencia se\~nalada
es de \textbf{amplitud de la evidencia reportada} sobre un mismo sistema,
no de calidad relativa del artefacto.

\section{Limitaciones de esta revisi\'on}
\label{sec:limitaciones-trabajos-relacionados}

\begin{itemize}[nosep]
  \item Solo se incluyeron fuentes en espa\~nol e ingl\'es; podr\'ian existir
    estudios relevantes en otros idiomas no cubiertos.
  \item El bloque de Jaime no document\'o el conteo de candidatos evaluados
    y excluidos (solo el resultado final de 4 incluidos), a diferencia de
    los bloques de Zaida (15 candidatos, 12 exclusiones documentadas) y de
    Fred (13 candidatos por DOI, 2 descartados con motivo individual);
    este gap se se\~nala expl\'icitamente con l\'inea punteada en el
    diagrama de flujo regenerado (figura~\ref{fig:prisma-flow}), \'unico
    bloque que a\'un no est\'a completamente cuantificado.
  \item \textbf{Inconsistencia num\'erica del bloque de Fred (detectada en
    auditor\'ia cruzada de Jaime el 2026-09-04, corregida el mismo d\'ia
    con la trazabilidad que Fred confirm\'o).} Una versi\'on anterior de
    este cap\'itulo y del diagrama de flujo reportaba, para el bloque
    Fred, ``13 candidatos $\to$ 2 excluidos $\to$ 3 incluidos'', dejando 8
    candidatos sin categor\'ia expl\'icita. La trazabilidad real es ``13
    identificados $\to$ 2 descartados $\to$ 11 retenidos $\to$ 3 trabajos
    relacionados + 8 de apoyo metodol\'ogico/rendimiento (no excluidos)'',
    ya reflejada en el p\'arrafo de trazabilidad de Fred arriba y en la
    figura~\ref{fig:prisma-flow}.
  \item La selecci\'on no tuvo un segundo revisor independiente por
    estudio: cada bloque fue filtrado por una sola persona sobre su propio
    tema, sin validaci\'on cruzada entre integrantes --- recurso real de un
    equipo de tres personas en un Proyecto Fin de Curso, declarado
    expl\'icitamente en vez de omitirse. El bloque de Fred aplic\'o,
    adem\'as, un filtro temporal expl\'icito (2020--2026) que los otros dos
    bloques no declararon como criterio formal --- diferencia de criterio
    entre bloques declarada, no oculta.
  \item Con el bloque de Fred, el dominio veterinario exacto de BIOPET pasa
    de un estudio de contraste a tres, pero ninguno combina el contexto
    ecuatoriano con una evaluaci\'on multim\'etrica del sistema: esa
    combinaci\'on espec\'ifica sigue subrepresentada en la literatura
    acad\'emica disponible, no por sesgo de b\'usqueda sino por escasez
    real, hallazgo reportado independientemente por los tres bloques de
    revisi\'on.
\end{itemize}
